\scanpage{46}

\mspage{- 176 -}\hfill IX

un homomorphisme canonique

\[
\tag*{(13.4.2)}
\mathcal{E}=\underline{\unsure{R|\mathbb{R}}}\otimes_{\mathbb{Z}}\underline{O}_S\longrightarrow\underline{E}\quad ,
\]

dont le noyau sera noté $\underline{F}(A)$ ou simplement $\underline{F}$, de sorte qu'on a une suite exacte

\[
\tag*{(13.4.3)}
0\longrightarrow\underline{F}\longrightarrow\mathcal{E}\longrightarrow\underline{E}\quad .
\]

Le cas le plus intéressant est celui où l'homomorphisme (13.4.2) est un épimorphisme, i.e. où pour tout $s\in S$, $R_s$ engendre le vectoriel complexe $E_s$. Dans ce cas on trouve donc une suite exacte

\[
\tag*{(13.4.4)}
0\longrightarrow\underline{F}\longrightarrow\mathcal{E}\longrightarrow\underline{E}\longrightarrow0\quad ,
\]

et si $R$ est localement constant, c'est là une suite exacte courte de Modules localement libres sur $S$. La connaissance du couple $(E,R)$ est alors équivalente à celle du couple $(R,\underline{F})$, où $R$ est un groupe analytique localement constant sur $S$ à fibres des $\mathbb{Z}$-Modules libres de type fini, et $\underline{F}$ est un sous-Module localement facteur direct de $\mathcal{E}=\underline{R}\otimes_{\mathbb{Z}}\underline{O}_S$, tel que l'on ait de plus la relation

\[
\tag*{(13.4.5)}
\underline{F}\cap\underline{R}=0\quad .
\]

Du point de vue des structures de Hodge, il y a lieu de regarder $\mathcal{E}$ comme filtré par la "\uline{filtration de Hodge}"

\[
\tag*{(13.4.6)}
\mathrm{Fil}^{1}(\mathcal{E})=\mathcal{E}\ ,\quad \mathrm{Fil}^{0}(\mathcal{E})=\underline{F}\ ,\quad \mathrm{Fil}^{-1}(\mathcal{E})=0\quad .
\]

Notons que la suite exacte (13.4.3) dépend encore fonctoriellement du groupe analytique $A$ sur $S$, et qu'elle commute aux changements de base quelconques.

\handwritten{\mspage{494}}

\mspage{494}

% archive sheet number: 45
